% Generated from locale/tr/content/first-order-logic/axiomatic-deduction/provability.tex; do not hand-edit.


\iftag{FOL}
      {\olfileid[tr]{fol}{axd}{prv-legacy}}
      {\olfileid[tr]{pl}{axd}{prv-legacy}}

\olsection{\usetoken{S}{derivability} Özellikleri}

\begin{prop}[Tekdüzelik] $\Gamma \subseteq \Delta$ ve
$\Gamma \Proves !A$ ise $\Delta \Proves !A$ da olur. \end{prop}

\begin{proof} Her sonlu $\Gamma_0 \subseteq \Gamma$, $\Delta$'nın da sonlu
bir altkümesidir; dolayısıyla $\Gamma_0 \Proves !A$ için !!a{derivation},
$\Delta \Proves !A$ olduğunu da gösterir. \end{proof}

\begin{prop}\ollabel{prop:provability} $\Gamma \Proves !A$ ancak ve ancak
$\Gamma\cup \{\lnot!A \}$ tutarsızdır. \end{prop}

\begin{prop} \begin{enumerate}
\item \ollabel{prop:provability-contr} $\Gamma \Proves !A$ ve
$\Gamma \cup \{ !A\} \Proves \lfalse$ ise $\Gamma$ tutarsızdır.

\item \ollabel{prop:provability-lnot} $\Gamma \cup \{!A\}
\Proves \lfalse$ ise $\Gamma \Proves \lnot !A$.

\item \ollabel{prop:provability-exhaustive} $\Gamma \cup \{!A\}
\Proves \lfalse$ ve $\Gamma \cup \{\lnot !A\} \Proves
\lfalse$ ise $\Gamma \Proves \lfalse$.

\tagitem{prvOr}{\ollabel{prop:provability-lor-left}
$\Gamma \cup \{!A\} \Proves \lfalse$ ve $\Gamma \cup \{!B\}
\Proves \lfalse$ ise $\Gamma \cup \{!A \lor !B\} \Proves \lfalse$.}{}

\tagitem{prvOr}{\ollabel{prop:provability-lor-right}
$\Gamma \Proves !A$ ya da $\Gamma \Proves !B$ ise
$\Gamma \Proves !A \lor !B$.}{}

\tagitem{prvAnd}{\ollabel{prop:provability-land-left}
$\Gamma \Proves !A \land !B$ ise $\Gamma \Proves !A$ ve
$\Gamma \Proves !B$.}{}

\tagitem{prvAnd}{\ollabel{prop:provability-land-right}
$\Gamma \Proves !A$ ve $\Gamma \Proves !B$ ise
$\Gamma \Proves !A \land !B$.}{}

\tagitem{prvIf}{\ollabel{prop:provability-mp}
$\Gamma \Proves !A$ ve $\Gamma \Proves !A \lif !B$ ise
$\Gamma \Proves !B$.}{}

\tagitem{prvIf}{\ollabel{prop:provability-lif}
$\Gamma \Proves \lnot !A$ ya da $\Gamma \Proves !B$ ise
$\Gamma \Proves !A \lif !B$.}{} \end{enumerate} \end{prop}

\begin{proof}
\begin{enumerate}
\item Sonlu $\Gamma_0,\Gamma_1\subseteq\Gamma$ için
$\Gamma_0 \cup \{!A\} \Proves \lfalse$ ile
$\Gamma_1 \Proves !A$ olduğunu varsayın.

\begin{derivation}
1. & $\Gamma_0 \cup \{!A\} \Proves \lfalse$ & varsayım \\
2. & $\Gamma_1 \Proves !A$ & varsayım\\
3. & $\Gamma_0 \cup \Gamma_1 \Proves \lfalse$ & kesme \\
\end{derivation}

$\Gamma_0 \subseteq \Gamma$ ve $\Gamma_1 \subseteq \Gamma$ olduğundan
$\Gamma_0 \cup \Gamma_1 \subseteq \Gamma$; dolayısıyla
$\Gamma \Proves \lfalse$.

\item Sonlu bir $\Gamma_0\subseteq\Gamma$ için
$\Gamma_0 \cup\{!A\} \Proves \lfalse$ olduğunu varsayın.

\begin{derivation}
1. & $\Gamma_0 \cup\{!A\} \Proves \lfalse$ & varsayım \\
2. & $\Gamma_0 \Proves !A \lif \lfalse$ & Tümdengelim Teoremi\\
3. & $\Gamma_0 \Proves \lnot !A$ & $\lnot !A$ tanımı \\
\end{derivation}

Tekdüzelik uyarınca $\Gamma \Proves \lnot !A$.

\item Sonlu $\Gamma_0,\Gamma_1\subseteq\Gamma$ için
$\Gamma_0 \cup \{ !A \} \Proves \lfalse$ ve
$\Gamma_1 \cup \{\lnot !A \} \Proves \lfalse$ olduğunu varsayın.

\begin{derivation}
1. & $\Gamma_0 \cup\{!A\} \Proves \lfalse$ & varsayım \\
2. & $\Gamma_0 \Proves !A \lif \lfalse$ & Tümdengelim Teoremi \\
3. & $\Gamma_1 \cup \{\lnot !A \} \Proves \lfalse$ & varsayım\\
4. & $\Gamma_0 \Proves ( !A \lif \lfalse) \lif \lnot !A$ & Ax0\\
5. & $\Gamma_0 \Proves \lnot !A$ & MP 2, 4\\
6. & $\Gamma_0 \cup \Gamma_1 \Proves \lfalse$ & Kesme 3, 5 \\
\end{derivation}

$\Gamma_0 \cup \Gamma_1 \subseteq \Gamma$ olduğundan
$\Gamma \Proves \lfalse$.

% prop:provability-lor-left
\tagitem{defOr}{}{
\iftag{probOr}{Alıştırma.}{Alıştırma.}}

% prop:provability-lor-right
\tagitem{defOr}{}{ \iftag{probOr}{Alıştırma.}{
$\Gamma_0 \Proves !A$ olduğunu varsayın.

\begin{derivation}
1. & $\Gamma_0 \Proves !A$ & varsayım \\
2. & $\Gamma_0 \Proves !A \lif (!A \lor !B)$ & Ax0\\
3. & $\Gamma_0 \Proves !A \lor !B$ & MP 1, 2 \\
\end{derivation}

Dolayısıyla $\Gamma \Proves !A \lor !B$. $\Gamma \Proves !B$ durumu da
benzerdir.}}

% prop:provability-land-left
\tagitem{defAnd}{}{ \iftag{probAnd}{Alıştırma}{
$\Gamma_0 \Proves !A \land !B$ olduğunu varsayın.

\begin{derivation}
1. & $\Gamma_0 \Proves !A \land !B$ & varsayım \\
2. & $\Gamma_0 \Proves (!A\land !B) \lif !A $ & Ax0\\
3. & $\Gamma_0 \Proves !A$ & MP 1, 2 \\
\end{derivation}

O halde $\Gamma \Proves !A$. Benzer bir !!{derivation},
$\Gamma \Proves !B$ olduğunu gösterir.}}

% prop:provability-land-right
\tagitem{defAnd}{}{\iftag{probAnd}{Alıştırma.}{Alıştırma.}}

% prop:provability-mp
\tagitem{defIf}{}{ \iftag{probIf}{Alıştırma.}{
$\Gamma_0 \Proves !A$ ve $\Gamma_1 \Proves !A \lif !B$ olsun.

\begin{derivation}
1. & $\Gamma_0 \Proves !A$ & varsayım \\
2. & $\Gamma_1 \Proves !A \lif !B $ & varsayım\\
3. & $\Gamma_0\cup\Gamma_1 \Proves !B$ & MP 1, 2\\
\end{derivation}

$\Gamma_0 \cup \Gamma_1 \subseteq \Gamma$ olduğundan
$\Gamma \Proves !B$.}}

% prop:provability-lif
\tagitem{defIf}{}{\iftag{probIf}{Alıştırma.}{Önce
$\Gamma \Proves \lnot !A$ olsun.

\begin{derivation}
1. & $\Gamma_0 \Proves \lnot !A$ & varsayım \\
2. & $\Gamma_0 \Proves \lnot !A \lif (!A \lif !B) $ & Ax0\\
3. & $\Gamma_0 \Proves !A \lif !B$ & MP 1, 2\\
\end{derivation}

$\Gamma \Proves !B$ durumunun ispatı benzerdir:

\begin{derivation}
1. & $\Gamma_0 \Proves !B$ & varsayım \\
2. & $\Gamma_0 \Proves !B \lif (!A \lif !B) $ & Ax0\\
3. & $\Gamma_0 \Proves !A \lif !B$ & MP 1, 2\\
\end{derivation}}}

\end{enumerate}
\end{proof}

\begin{probtag}{probOr,probAnd,probIf}
\olref[fol][axd][prv-legacy]{prop:provability} önermesinin ispatını tamamlayın.
\end{probtag}

\begin{thm}[Genelleme] \ollabel{thm:Generalization}
$\Gamma \Proves !A$ ve $x$, $\Gamma$ içindeki hiçbir !!{formula}{}de serbest
değilse $\Gamma \Proves \lforall[x][!A]$.
\end{thm}

\begin{proof} $\mathsf{Thm}_\Gamma$ üzerinde tümevarımla: $!A$ bir aksiyomsa
$\lforall[x][!A]$ da aksiyomdur. $!A\in\Gamma$ ise $x$, $!A$ içinde serbest
değildir; dolayısıyla $!A \lif \lforall[x][!A]$ bir aksiyomdur
(\textbf{Ax3}) ve MP ile $\Gamma \Proves \lforall[x][!A]$. Şimdi $!A$'nın,
$\Gamma \Proves !B$ ve $\Gamma \Proves !B \lif !A$ olduğu için MP ile
çıktığını varsayın. Tümevarım varsayımına göre
$\Gamma \Proves \lforall[x][!B]$ ve
$\Gamma \Proves \lforall[x][(!B \lif !A)]$. \textbf{Ax2} uyarınca ayrıca
\[
\Gamma \Proves \lforall[x][(!B \lif !A)] \lif
(\lforall[x][!B] \lif \lforall[x][!A]);
\]
MP'nin iki uygulaması istendiği gibi
$\Gamma \Proves \lforall[x][!A]$ verir. \end{proof}

\begin{thm}\ollabel{thm:WeakGen} (\emph{Sabitler Üzerinde Zayıf Genelleme})
$\Gamma \Proves !A$ ve $c$, $\Gamma$ içinde bulunmayan !!a{constant} ise
$!A$ içinde bulunmayan bir $x$ değişkeni için
$\Gamma \Proves \lforall[x][\Subst{!A}{x}{c}]$ ve ispat $c$ içermez.
\end{thm}

\begin{proof} $!A_1,\ldots,!A_n$, $\Gamma$'dan $!A$'nın bir ispatı ve
$!A=!A_n$ olsun. $!A_1,\ldots,!A_n$ içinde bulunmayan bir $x$ değişkeni seçip
şu yeni diziyi ele alın:
$\Subst{!A_1}{x}{c},\ldots,\Subst{!A_n}{x}{c}$. Bu, $\Gamma$'dan
$\Subst{!A}{x}{c}$'nin bir ispatıdır. Gerçekten her $i=1,\ldots,n$ için:
\begin{itemize}
\item $!A_i$ bir aksiyomsa $\Subst{!A_i}{x}{c}$ de aksiyomdur;
\item $!A_i \in \Gamma$ ise $c$, $\Gamma$ içinde bulunmadığından
$\Subst{!A_i}{x}{c} = !A_i \in \Gamma$;
\item $!A_i$, $!A_j$ ile $!A_j \lif !A_i$'dan elde edilmişse
$\Subst{!A_i}{x}{c}$, $\Subst{!A_j}{x}{c}$ ile
$\Subst{(!A_j \lif !A_i)}{x}{c} = \Subst{!A_j}{x}{c} \lif
\Subst{!A_i}{x}{c}$'dan MP ile çıkar.
\end{itemize}
Yeni dizide !!{constant} $c$'nin artık bulunmadığı açıktır. $\Gamma'$,
$\Subst{!A}{x}{c}$'nin !!{derivation}{}inde görünen, $\Gamma$ içindeki
!!{formula}s{}den oluşsun. O zaman $x$, $\Gamma'$ içindeki hiçbir !!{formula}{}de
serbest değildir ve $\Gamma' \Proves \Subst{!A}{x}{c}$. Genelleme
(\olref{thm:Generalization}) uyarınca
$\Gamma' \Proves \lforall[x][\Subst{!A}{x}{c}]$; tekdüzelik uyarınca
istendiği gibi $\Gamma \Proves \lforall[x][\Subst{!A}{x}{c}]$.
\end{proof}

\begin{explain}
Amacımız \olref{thm:WeakGen} içindeki $x$'in $!A$ içinde hiç bulunmaması
koşulunu, $x$'in $!A$ içinde serbest olmaması ve $!A$ içinde $c$ için
!!{free for} olması biçimindeki daha zayıf koşullarla değiştirmektir. Bunun
bağlı !!{variable} değiştirilerek yapılabileceği açıktır; önce bunu
ispatlayacağız.
\end{explain}

\begin{lem}
\ollabel{lem:free-for}
$x$, $!A$ içinde $c$ için serbest ve $y$, $!A$ içinde serbest değilse $x$,
$\Subst{!A}{y}{c}$ içinde $y$ için !!{free for}{}dir.
\end{lem}

\begin{proof} $x$, $\Subst{!A}{y}{c}$ içinde $y$ için !!{free for} değilse
$\Subst{!A}{y}{c}$ içindeki serbest bir $y$ oluşumu bir $\lforall[x]$
niceleyicisinin kapsamında kalır. Varsayım gereği $y$, $!A$ içinde serbest
olmadığından bütün bu oluşumlar $\Subst{}{y}{c}$ yerine koymasından gelir.
O halde $c$'nin bir oluşumu $\lforall[x]$ kapsamında kalır ve $x$, $!A$
içinde $c$ için !!{free for} değildir; bu bir çelişkidir.
\end{proof}

\begin{lem}[Bağlı Değişken Değiştirme]
\ollabel{lem:ChangeBdVar}
$x$ ile $y$, $!A$ içinde serbest değil ve ikisi de $!A$ içinde $c$ için
!!{free for} ise
$\Proves \lforall[x][\Subst{!A}{x}{c}] \equiv
\lforall[y][\Subst{!A}{y}{c}]$ (ve ispat $c$ içermez).
\end{lem}

\begin{proof}
Varsayımlara göre $y$, $!A$ içinde $c$ için !!{free for} ve $x$, $!A$
içinde serbest değildir. Önceki \olref{lem:free-for} uyarınca $y$,
$\Subst{!A}{x}{c}$ içinde $x$ için !!{free for}{}dir. Dolayısıyla
$\lforall[x][\Subst{!A}{x}{c} \lif \Subst{!A}{x}{c} \Subst{}{y}{x}]$
bir aksiyomdur (\textbf{Ax1}). $x$, $!A$ içinde serbest olmadığından
$\Subst{!A}{x}{c} \Subst{}{y}{x} = \Subst{!A}{y}{c}$ ve dolayısıyla
\[
\Proves \lforall[x][\Subst{!A}{x}{c}] \lif \Subst{!A}{y}{c}.
\]
Tümdengelim Teoremiyle
$\lforall[x][\Subst{!A}{x}{c}] \Proves \Subst{!A}{y}{c}$. $y$, $!A$
içinde serbest olmadığından $\lforall[x][\Subst{!A}{x}{c}]$ içinde de
serbest değildir. Genelleme uyarınca
$\lforall[x][\Subst{!A}{x}{c}] \Proves
\lforall[y][\Subst{!A}{y}{c}]$; Tümdengelim Teoremi bir kez daha
$\Proves \lforall[x][\Subst{!A}{x}{c}] \lif
\lforall[y][\Subst{!A}{y}{c}]$ verir. Ters yönün ispatı bütünüyle
simetriktir; sonuç Taut ile çıkar.
\end{proof}

\begin{thm}[Sabitler Üzerinde Güçlü Genelleme]
\ollabel{thm:StrongGen}
$\Gamma \Proves !A$, !!{constant} $c$, $\Gamma$ içinde bulunmuyor, $x$ ise
$!A$ içinde serbest olmayıp $c$ için !!{free for} ise
$\Gamma \Proves \lforall[x][\Subst{!A}{x}{c}]$.
\end{thm}

\begin{proof}
$\Gamma \Proves !A$ ve $c$, $\Gamma$ içinde bulunmadığından Zayıf Genelleme
uyarınca $!A$ içinde bulunmayan !!a{variable} $y$ için
$\Gamma \Proves \lforall[y][\Subst{!A}{y}{c}]$. $y$, $!A$ içinde hiç
bulunmadığından serbest değildir ve $c$ için !!{free for}{}dir. Varsayım gereği
$x$ de $!A$ içinde serbest değil ve $c$ için !!{free for} olduğundan bağlı
!!{variable} değiştirme koşulları sağlanır. Dolayısıyla
$\Proves \lforall[x][\Subst{!A}{x}{c}] \equiv
\lforall[y][\Subst{!A}{y}{c}]$ ve buradan
$\Gamma \Proves \lforall[x][\Subst{!A}{x}{c}]$ çıkar.
\end{proof}

\begin{thm}
\ollabel{thm:provability-quantifiers}
\begin{tagenumerate}{prvEx,prvAll}
\tagitem{prvEx}{$\Gamma \Proves !A(t)$ ise
$\Gamma \Proves \lexists[x][!A(x)]$.}{}
\tagitem{prvAll}{$\Gamma \Proves \lforall[x][!A(x)]$ ise
$\Gamma \Proves !A(t)$.}{}
\end{tagenumerate}
\end{thm}

\begin{proof}
\begin{tagenumerate}{prvEx,prvAll}
%% Varoluş için aksiyomlar gerekiyor
\tagitem{prvEx}{Alıştırma.}{}
\tagitem{prvAll}{$\Gamma_0 \Proves \lforall[x][!A(x)]$ olduğunu varsayın.

\begin{derivation}
1. & $\Gamma_0 \Proves \lforall[x][!A(x)]$ & varsayım \\
2. & $\Gamma_0 \Proves \lforall[x][!A(x)] \lif \Subst{!A}{t}{x}$ & Ax1 \\
3. & $\Gamma_0 \Proves \Subst{!A}{t}{x}$ & MP 1, 2 \\
\end{derivation}}{}

\end{tagenumerate}
\end{proof}

